\section{Discussion}
\label{sec:discussion}

\subsection{Limitations}
\label{sec:limitations}

\draft{The approach depends on our SpecTec-to-Lean translation, the implementation connection, and the properties the agent succeeds in proving.}
\draft{Translation-validation tests have passed, but they cannot establish semantic preservation for all SpecTec inputs.}
\draft{Its scope is also limited by the properties expressible in the supported SpecTec-like fragment and by the available generator and oracle derivation procedures.}
% TODO: State the concrete expressiveness limits and the effect of execution or generation bounds once the language is fixed.
\draft{Completeness concerns reachability of inputs, not their sampling frequency, and a valid property may still be too weak to find bugs.}

\subsection{Beyond Language Specifications}
\label{sec:beyond}

\draft{SpecTrum uses Consensus-SpecTec, a mechanized specification of the Ethereum consensus algorithm, to guide differential fuzzing of consensus clients~\cite{jeong26spectrum}.}
\draft{This provides a concrete setting beyond language specifications in which to explore our approach.}

\draft{The approach could apply beyond P4 where a specification is available in a theorem prover, the desired properties fit the supported fragment, and implementation observations can be related to it.}
\draft{Such an extension also requires support for deriving generators from the relevant premises.}
\draft{A checked decision procedure is also needed for each conclusion used as an oracle; applicability outside P4 remains to be evaluated.}
